\documentclass{article}
\usepackage{graphicx}
\usepackage{booktabs}
\usepackage{amsmath}
\usepackage{hyperref}
\usepackage{multirow}

\title{AnnotateBench: A Systematic Benchmark of Annotation Strategies Across NLP Task Types}
\author{[Intern Name], Spurthi Setty \\
Anote AI}
\date{}

\begin{document}

\maketitle

\begin{abstract}
\noindent
How much labeled data does a given NLP task actually need? The answer depends on both the task type and the annotation strategy used, yet no existing work systematically studies this relationship across task types at scale. We introduce \textbf{AnnotateBench}, a benchmark of 750 experimental configurations spanning 5 annotation strategies, 10 datasets across 5 NLP task types, and 5 annotation budget levels. We fit learning curves for each configuration and derive cost-accuracy Pareto frontiers. Our key finding is that no single annotation strategy dominates across all task types: active learning strategies excel at classification and NER but underperform on complex generation tasks where diversity-based sampling outperforms uncertainty-based methods. We release all experimental code, fitted curves, and a practitioner decision tree.
\end{abstract}

\section{Introduction}
% Annotation cost is the primary bottleneck in supervised NLP
% Prior work studies single strategies on single tasks
% This paper: systematic cross-task benchmark

\section{Related Work}
\subsection{Active Learning for NLP}
% BADGE, BERT-based uncertainty sampling, diversity-based
\subsection{LLM-as-Annotator}
% Gilardi 2023, Ziegler 2022
\subsection{Data-Centric AI}
% DataPerf, Data Maps, annotation efficiency work

\section{AnnotateBench Design}
\subsection{Task Types and Datasets}
\begin{table}[h!]
\centering
\caption{AnnotateBench Datasets}
\begin{tabular}{lll}
\toprule
\textbf{Task Type} & \textbf{Dataset 1} & \textbf{Dataset 2} \\
\midrule
Text Classification & SST-2 & AG News \\
Named Entity Recognition & CoNLL-2003 & OntoNotes \\
Extractive QA & SQuAD 2.0 & Natural Questions \\
Abstractive Summarization & CNN/DailyMail & XSum \\
Instruction Tuning & Alpaca & FLAN collection subset \\
\bottomrule
\end{tabular}
\end{table}

\subsection{Annotation Strategies}
% (1) Random Sampling, (2) Uncertainty Sampling, (3) Diversity Sampling (CoreSet),
% (4) BADGE, (5) LLM-as-Annotator (GPT-4o zero-shot)

\subsection{Budget Levels}
% 10 / 50 / 100 / 500 / 1,000 labeled examples

\subsection{Models}
% BERT-base/large (classification, NER), T5/FLAN-T5-base (QA, summarization),
% Llama 3 8B SFT (instruction tuning)

\section{Experimental Setup}
% 5 strategies × 10 datasets × 5 budgets × 3 models = ~750 configs
% Learning curve fitting: power law y = a * x^b + c
% Pareto frontier computation

\section{Results}
% Figure: Learning curves by strategy across task types
% Table: Best strategy per task type at each budget level
% Figure: Cost-accuracy Pareto frontiers

\subsection{No Single Strategy Dominates}
\subsection{Task Complexity Predicts Annotation Sufficiency}
\subsection{LLM-as-Annotator: When It Works and When It Fails}

\section{Analysis}
\subsection{Fitted Learning Curve Parameters}
\subsection{Practitioner Decision Tree}

\section{Conclusion}

\section*{References}
\small
\begin{enumerate}
    \item Ash et al. BADGE: Batch Active learning by Diverse Gradient Embeddings, 2020. ICLR 2020.
    \item Gilardi et al. ChatGPT Outperforms Crowd Workers for Text-Annotation Tasks, 2023. PNAS.
    \item Sener \& Savarese. Active Learning for Convolutional Neural Networks: A Core-Set Approach, 2018. ICLR 2018.
    % Add more references
\end{enumerate}

\end{document}
